\chapter{Научный контракт и проверяемые гипотезы}

Контракт разводит ноксический стимул, ноцицептивную реакцию, защитное поведение,
самоотчёт боли, ECAP-рекрутирование и клинический ответ на SCS.

\section{Гипотезы}
H1 относится только к воспроизводимости симуляции; H2 --- к приросту на независимых
человеческих данных; H3 --- к прогнозу заранее заданного SCS-исхода на новых пациентах.

\section{STOP-критерии}
Для H1--H3 заранее определены условия остановки, резервный результат и запрещённый вывод.
Отрицательный результат сужает работу, но не переименовывает её цель.
